\subsection{\bsq{if}, \bsq{if-else}, \bsq{if-else~if-else} Statements}\label{sec:IfStatements}

We already discussed the use of blanks, round brackets, curly braces and linefeeds in \tqfullvref{sec:WhiteSpace}.

\keeplisting{11}
The base form for the \lstinline|if| statement is like this:
\begin{lstlisting}
if( <condition>)
{
    // <condition> is true
    <statements>;
}
else
{
    // <condition> is false
    <statements>;
}
\end{lstlisting}

\keeplisting{7}
If there is nothing to do in case the condition results as false, the \lstinline|else| clause must be omitted:
\begin{lstlisting}
if( <condition>)
{
    // <condition> is true
    <statements>;
}
\end{lstlisting}

\keeplisting{8}
Negate the condition if there is nothing to do when the original condition is true; do not let the first statement block empty:
\begin{lstlisting}
// RECOMMENDED!!
if( !<condition>)
{
    // <condition> is false; that means !<condition> is true
    <statements>;
}
\end{lstlisting}

\keeplisting{12}
Instead of:
\begin{lstlisting}
// AVOID!!
if( <condition>)
{
    // <condition> is true
    // Nothing to do here!
}
else
{
    // <condition> is false
    <statements>;
}
\end{lstlisting}

\keeplisting{3}
If there is no \lstinline|else| clause, and only a single statement, you can write the \lstinline|if| statement like this:
\begin{lstlisting}
if( <condition> ) <statements>;
\end{lstlisting}

The curly braces are mandatory if it is not possible to write the whole \lstinline|if| statement within a single line.

\keeplisting{22}
In case there are more conditions to check, but a \lstinline|switch| statement cannot be used due to the data types involved or the logic for the conditions, use this form:
\begin{lstlisting}
if( <condition1>)
{
    <statements>;
}
else if( <condition2>)
{
    <statements>;
}
else if( <condition3>)
{
    <statements>;
}
else if( <condition#>)
{
    <statements>;
}
else
{
    <statements>;
}
\end{lstlisting}

Of course you can repeat the \lstinline|else if| as often as you need, but at some point the readability suffers seriously, and you should think about another logic. Also the final \lstinline|else| clause is not mandatory.

\keeplisting{33}
It is not recommended to write \verb#if-else# like this, although it seems to be more conformant to the basic rule:
\begin{lstlisting}
// NOT RECOMMENDED!!
if( <condition1>)
{
    <statements>;
}
else 
{
    if( <condition2>)
    {
        <statements>;
    }
    else 
    {
        if( <condition3>)
        {
            <statements>;
        }
        else 
        {
            if( <condition#>)
            {
                <statements>;
            }
            else
            {
                <statements>;
            }
        }
    }
}        
\end{lstlisting}

The increasing indendation makes this even worse to read than the recommended variant, particularly when the code lines get longer.

\subsubsection{Principles}
\begin{enumerate}[label=P\arabic*.]
\item{Curly braces are mandatory when the \lstinline|if| statement has more than one line. They are always required for the \lstinline|else| clause.}
\item{No empty \lstinline|if| clause; negate the condition if necessary.}
\item{No blanks between the \lstinline|if|\index{if} keyword and the opening round bracket. 

Refer to \tqfullvref{sec:ParameterParenthesis} on how to place the white space.}
\item{The opening curly brace for the body statement of the \lstinline|if| and the optional \lstinline|else| clause is on the same indentation level as the keyword itself. In case of the \verb#if-else# construct, on the level of the \lstinline|else|.}
\item{The statements itself are indented one level relative to the curly braces.}
\end{enumerate}

Please refer also to chapter \tqvref{sec:TheTernaryOperator} that discusses the ternary \verb#?# operator.

%==============================================================================
% End of file!!
%==============================================================================
